\section{Методология исследования}

В данном разделе описывается математический аппарат и алгоритмическая последовательность действий, использованных для проведения статистического анализа и оценки характеристик брендов...

\subsection{Метод идеальной точки}

Для количественной оценки потребительских предпочтений в работе используется \textit{модель идеальной точки}. Данный метод позволяет оценить степень близости реальных характеристик товара к ожиданиям потребителя. Расчет итогового балла бренда ($A_b$) производится по формуле:

\begin{equation}
    \label{eq:integral_score}
    A_{b} = \sum_{i=1}^{m} W_{i} \cdot |I_{i} - X_{i}|  ,
\end{equation}
где 
\begin{description}
    \item[$m$] --- количество параметров, 
    \item[$W_i$] --- коэффициент важности $i$-го параметра, 
    \item[$I_i$] --- значение $i$-го параметра для «идеального» продукта, 
    \item[$X_i$] --- фактическое значение $i$-го параметра оцениваемого бренда.
\end{description}

Чем ниже полученное значение $A_b$, тем выше лояльность потребителя к бренду, так как его характеристики ближе к идеальным.

\subsection{Формирование выборки и методы отбора}

Центральной задачей работы является сравнение эффективности различных методов формирования выборки из генеральной совокупности (данных опроса респондентов). Были реализованы следующие алгоритмы:

\begin{enumerate}
    \item \textbf{Собственно-случайный отбор (Simple Random Sampling)}\\
    Базовый метод статистического исследования, основанный на принципе равновероятности включения любого элемента генеральной совокупности в выборку. Метод минимизирует систематическую ошибку исследователя и используется в качестве контрольного уровня для оценки статистической эффективности других алгоритмов. В данной работе реализована модель \textit{бесповторного} и \textit{повторного} отбора.

    \item \textbf{Механический отбор (Mechanical Sampling)} \\
    Систематический метод, предполагающий отбор единиц из упорядоченного массива данных через равные интервалы. Метод эффективен при отсутствии периодических закономерностей в исходном массиве, которые могли бы совпасть с шагом отбора и исказить результаты. Такой метод отбора подразумевает только \textit{бесповторный} способ.
    
    \item \textbf{Типический отбор (Stratified Sampling)} \\
    Метод, основанный на предварительном разделении (стратификации) неоднородной совокупности на качественно однородные страты по заданному признаку\footnote{В работе применялась поло-возрастная стратификация.}. В рамках исследования выборка формируется путем случайного отбора из каждой страты пропорционально их доле в генеральной совокупности. Данный подход позволяет существенно снизить стандартную ошибку средней за счет исключения межгрупповой дисперсии из общей ошибки выборки. В данной работе реализована модель \textit{бесповторного} и \textit{повторного} отбора.
    
    \item \textbf{Кластерный отбор (Cluster Sampling)} \\
    Метод формирования выборки, при котором объектом отбора выступают не отдельные респонденты, а их логические или структурные группы~---~кластеры\footnote{Деление на кластеры в работе происходило по географическому признаку.}. Выборка формируется путем случайного отбора самих серий, внутри которых проводится сплошное обследование. Использование данного метода в работе позволяет проанализировать устойчивость оценок брендов с учетом внутригрупповой корреляции и оценить влияние специфических характеристик конкретных групп данных на итоговый интегральный признак. В данной работе реализована модель \textit{бесповторного} и \textit{повторного} отбора.
\end{enumerate}

\newpage

\subsection{Оценка точности и надежности результатов}
Для того чтобы полученные результаты были \textit{репрезентативными} (т.е. адекватно воспроизводили структуру генеральной совокупности), в работе проводится расчет следующих показателей:

\begin{enumerate}
    \item \textbf{Выборочная средняя оценка} --- среднее арифметическое значений всех объектов в выборке, выступающее в роли точечной оценки генерального среднего значения:
    \begin{equation}
        \label{eq:mean}
        \bar{x} = \frac{1}{n} \sum_{i=1}^{n} x_i  ,
    \end{equation}
    где $x_i$ --- значение оценки $i$-го респондента, $n$ --- объем выборки.

    \item \textbf{Система показателей вариации} \\
    При разделении совокупности на $k$ групп (страт или кластеров) по определенному признаку, вариация признака анализируется через систему взаимосвязанных показателей:

    \begin{itemize}
        \item \textbf{Выборочная дисперсия ($\hat{\sigma}^2$)} --- характеризует вариацию признака во всей совокупности под влиянием всех факторов:
        \begin{equation}
            \label{eq:var}
            \hat{\sigma}^2 = \frac{\sum_{i=1}^{n} (x_i - \bar{x})^2}{n-1}  ,
        \end{equation}

        где 
        \begin{description}
            \item[$x_i$] --- $i$-ое значение из выборки, 
            \item[$\bar{x}$] --- среднее по всей выборке, 
            \item[$n$] --- размер выборки.
        \end{description}

        \clearpage
        
        \item \textbf{Внутригрупповая дисперсия ($\hat{\sigma_j}^2$)} --- отражает случайную вариацию в каждой отдельной $j$-й группе, вызванную факторами, не принятыми за основание группировки:
        \begin{equation}
            \label{eq:var_within}
            \hat{\sigma_j}^2 = \frac{\sum_{i=1}^{n_j} (x_{ij} - \bar{x}_j)^2}{n_j - 1},  j=\overline{1,k}
        \end{equation}

        где 
        \begin{description}
            \item[$x_{ij}$] --- $i$-ое значение из $j$-ой группы; 
            \item[$\bar{x}_j$] --- среднее по всей $j$-ой группе; 
            \item[$n_j$] --- размер $j$-ой группы.
        \end{description}

        \item \textbf{Средняя из внутригрупповых дисперсий ($\overline{\hat{\sigma}^2}$)} --- характеризует вариацию, обусловленную всеми факторами, за исключением фактора группировки:
        \begin{equation}
            \label{eq:var_within_mean}
            \overline{\hat{\sigma}^2} = \frac{\sum_{j=1}^{k} \hat{\sigma_j}^2 \cdot n_j}{\sum_{j=1}^{k} n_j}  ,
        \end{equation}

        где 
        \begin{description}
            \item[$s_j$] --- внутригрупповая дисперсия, 
            \item[$n_j$] --- размер $j$-ой группы.
        \end{description}
        
        \item \textbf{Межгрупповая дисперсия ($\delta^2$)} --- характеризует вариацию результативного признака, обусловленную исключительно влиянием фактора, положенного в основу группировки:
        \begin{equation}
            \label{eq:var_between}
            \delta^2 = \frac{\sum_{j=1}^{k} (\bar{x}_j - \bar{x})^2 \cdot n_j}{\sum_{j=1}^{k} n_j}  ,
        \end{equation}

        где 
        \begin{description}
            \item[$\bar{x}_j$] --- среднее $j$-ой группы, 
            \item[$\bar{x}$] --- среднее по всем группам, 
            \item[$n_j$] --- размер $j$-ой группы.
        \end{description}
        
    \end{itemize}

    Согласно \textit{правилу сложения дисперсий}, общая дисперсия равна:
    \begin{equation}
        \hat{\sigma}^2 = \overline{\hat{\sigma}^2} + \delta^2.
    \end{equation}

    \clearpage
    
    \item \textbf{Предельная ошибка выборки} --- максимально возможное отклонение выборочной средней от 
    генеральной для заданной доверительной вероятности. Рассчитывается по формуле:
    \begin{equation}
        \label{eq:delta_error}
        \Delta = t \cdot \sqrt{\frac{\hat{\sigma}_{sample}^2}{n}} \cdot \eta  ,
    \end{equation}
    
    где
    \begin{description}
        \item[$\hat{\sigma}_{sample}^2$] --- выборочная дисперсия метода,
        \item[$t$] --- квантиль распределения, определяемый уровнем
        значимости $\alpha$\textsuperscript{\ref{fn:alpha}},
        \item[$\eta$] --- поправка на конечный объем совокупности \eqref{eq:eta},
        \item[$n$] --- объемы выборки,
        \item[$N$] --- объем генеральной совокупности. 
    \end{description}

    \footnotetext{В работе уровень значимости принимается равным $\alpha = 0,05$, что соответствует 
    доверительной вероятности $P = 0,95$ \label{fn:alpha}}

    \begin{equation}
        \label{eq:eta}
        \eta = 
        \begin{cases} 
        \sqrt{1-\dfrac{n}{N}} & (\text{для бесповторного отбора}), \\
        1, & (\text{для повторного отбора}).
        \end{cases}
    \end{equation}

\end{enumerate}



\subsection{Оценка и обоснование объема выборки}

Для подтверждения \textit{репрезентативности} данных в работе проводится расчет \textit{минимально необходимого объема выборки} ($n_{min}$) с использованием итерационной схемы, согласно требованиям технического задания.



\hypertarget{algo_target}{\subsubsection*{Итерационный алгоритм расчета}}

Процесс вычисления $n_{min}$ разделен на три этапа, что позволяет
учесть форму распределения при малых выборках и конечность генеральной совокупности:

\begin{enumerate}
    \item \textbf{Расчет первого приближения ($n_0^*$)}.
    Используется квантиль стандартизированного нормального распределения $z$
    для заданного уровня значимости $\alpha$\textsuperscript{\ref{fn:alpha}}:

    \begin{equation}
        \label{eq:n_0}
        n_0^* = \frac{z^2 \cdot \hat{\sigma}_{sample}^2}{\Delta_{GP}^2}  ,
    \end{equation}

    где 
    \begin{description}
        \item[$\hat{\sigma}_{sample}^2$] --- соответствующий методу вид дисперсии, 
        \item[$\Delta_{GP}$] --- соответствующая методу предельная ошибка генеральной совокупности.\footnotemark
    \end{description}
    \footnotetext{Согласно техническому заданию, оценка производится по
    уполовиненной величине предельной ошибки соответствующего метода: $\Delta_{GP} = \frac {\Delta} {2}$ \label{fn:delta}.}
    
    \item \textbf{Уточнение через распределение Стьюдента ($n_1^*$)}.
    На основе полученного $n_0^*$ определяется число степеней свободы, и вычисляется
    уточненный коэффициент доверия $t_{new}$ с помощью распределения Стьюдента:
    \begin{equation}
        \label{eq:n_1}
        n_1^* = \frac{t_{new}^2 \cdot \hat{\sigma}_{sample}^2}{\Delta_{GP}^2}
    \end{equation}

    \item \textbf{Финальная корректировка ($n^*$)}.
    Для методов, предполагающих бесповторный отбор, производится пересчет 
    с учетом объема генеральной совокупности $N$:
    \begin{equation}
        n^* = 
        \begin{cases} 
            n_1^* & \quad (\text{для повторного способа}), \\
            \dfrac{N \cdot n_1^*}{N + n_1^*} & \quad (\text{для бесповторного способа}).
        \end{cases}
    \end{equation}

\noindent Итоговый \textit{минимальный объем выборки} принимается как $n_{min} = \lceil n^* \rceil$.

\end{enumerate}

\hypertarget{z-test}{\subsection{Методика проверки статистической несмещенности на основе Z-критерия}}

Для оценки качества работы алгоритмов формирования выборки в рамках имитационного моделирования
был применен \textit{статистический тест} для результатов.
Ключевым критерием точности метода в данной работе выступает свойство несмещенности оценки 
математического ожидания.

Поскольку параметры исходной базы данных\footnote{Принимается за генеральную совокупность в работе.}
известны полностью, для проверки гипотезы о равенстве выборочного 
среднего ($\bar{x}$) и истинного среднего ($\mu$) используется двухсторонний $Z$-тест. Применение 
данного критерия обосновано следующими факторами:

\begin{enumerate}[noitemsep]
    \item Известна \textbf{дисперсия генеральной совокупности} ($\sigma^2$).
    \item Ввиду \textit{центральной предельной теоремы} распределение средних значений, полученных в ходе достаточного большого числа итераций, стремится к \textbf{нормальному распределению} $\mathcal{N}(\mu, \frac{\sigma^2}{n})$ независимо от закона распределения исходных данных в мастер-таблице.
\end{enumerate}

Математическая модель проверки строится на тестировании системы гипотез:
\begin{equation}
\begin{cases}
H_0: \bar{x} = \mu \quad (\text{нулевая гипотеза}), \\
H_1: \bar{x} \neq \mu \quad (\text{альтернативная гипотеза}).
\end{cases}
\end{equation}

Для каждого метода и исследуемого признака рассчитывается эмпирическое значение $Z$-статистики:
\begin{equation}
    Z_{emp} = \frac{\bar{x} - \mu}{SE_{sim}}
\end{equation}

Стандартная ошибка $SE_{sim}$ в контексте имитационного моделирования рассчитывается по формуле:
\begin{equation}
    SE_{sim} = \sqrt {\frac{\hat{\sigma}^2}{n \cdot k}} \cdot \eta
\end{equation}
где 
\begin{description}
    \item[$\hat{\sigma}^2$] --- несмещенная оценка дисперсии метода, 
    \item[$n$] --- объем единичной выборки, 
    \item[$k$] --- количество итераций моделирования, 
    \item[$\eta$] --- поправочный коэффициент, зависящий от типа отбора \eqref{eq:eta}.
\end{description}

\bigskip

Нулевая гипотеза отвергается, если расчетное значение $p\text{-}\text{\textsmaller{value}}$ оказывается меньше $\alpha$\textsuperscript{\ref{fn:alpha}}:
\begin{equation}
    p\text{-}\text{\textsmaller{value}} = 2 \cdot \left(1 - \Phi(|Z_{emp}|)\right)
\end{equation}
где $\Phi$ --- функция Лапласа.

\begin{equation}
    \Phi(x) = \frac{1}{\sqrt{2\pi}} \int_{0}^{x} e^{-\frac{t^2}{2}} dt.
\end{equation}

В случае отклонения гипотезы $H_0$ метод признается смещенным для данной конфигурации параметров.